\section{Compact finiteness from stable transfer images}
\label{sec:compact}

We now prove the compact part of Theorem~\ref{thm:coefficient-finiteness}.
The support comparison supplies a finite stable image for transfer.
The argument below uses the translation representations to turn this
stable-image statement into countability for the localized coefficients,
then uses compactness to obtain finiteness for the compact ones.
The bounded-pole finiteness assertion remains the task of
Section~\ref{sec:boundary}.

For $P$ as in Lemma~\ref{lem:compact:stable-lattices},
Corollary~\ref{support:finite-stable-image} gives
\begin{equation}\label{eq:compact:finite-stable-part}
 \dim_k\bigcap_{a\geq0}S^aH^i(P_n,P)<\infty.
\end{equation}

\subsection{Pole strips and the compact algebra lemma}

Assume the earlier compact assertions in
\eqref{eq:compact:induction-order}.  Every strip between two fixed pole
bounds for the lattices above has finite $P_n$-cohomology.  Indeed, write
$x$ as a unit times the product of the nonzero label coordinates.
Successively dividing out one factor filters a fixed strip by finitely
many restrictions to a label hyperplane, with its conormal powers.
That hyperplane is a starting-height $t+1$ basis disc after the finite
root change in Corollary~\ref{rings:boundary-labels}.  Continuing through
intersections gives only the admissible boundary coefficients in the
induction hypothesis.  Power-series truncations have continuous module
sections, so these short exact sequences remain exact on continuous
cochain terms.  At $t=n$ this is simply finite-coefficient group
cohomology.

Write, at a fixed frame level,
\[
             M^i=H^i(P_n,P),\qquad X^i=H^i(P_n,B_U).
\]
The map $M^i\to X^i$ has countable-dimensional kernel and cokernel.
To see this, express the localized coefficient as the union of
$x^{-c}L$ for integral $c$.  The finite resolution commutes with this
union, and the differences from $P$ are the strips just considered.
There are countably many bounds and each strip has finite cohomology.
Moreover the kernel is locally $S$-nilpotent.  If a cocycle in $P$
becomes a boundary after localization, choose a primitive at one pole
bound.  A high iterate of $S$ carries that primitive into $P$, and
annihilates the original cohomology class.

\begin{lemma}\label{lem:compact:operator}
Let $M$ be a compact $k$-vector space and $S$ a continuous $k$-linear
endomorphism.  Suppose $M/SM$ and $I=\bigcap_{a\geq0}S^aM$ are finite.
Then $S(I)=I$, $S$ is invertible on $I$, and $M/I$ is a finitely
generated $k[[z]]$-module with $z$ acting as $S$.  In particular $S$
has finite kernel and cokernel on $M$.

If $f:M\to X$ commutes with $S$ and has locally $S$-nilpotent kernel,
then $\ker f$ is finite and
\begin{equation}\label{eq:compact:finite-intersection}
 \bigcap_{a\geq0}\operatorname{im}(S^aM\longrightarrow X)=f(I)
\end{equation}
is finite.
\end{lemma}

\begin{proof}
For $y\in I$, the nonempty compact sets
$S^{-1}(y)\cap S^aM$ are nested.  Their intersection gives a preimage
in $I$, proving $S(I)=I$; finiteness makes the restriction invertible.
The images $S^a(M/I)$ shrink uniformly to zero.  Otherwise a compact
closed set outside a neighborhood of zero would meet every one of
these nested images and hence their intersection, which is zero.

Choose finitely many lifts of a basis of $(M/I)/S(M/I)$.  Repeatedly
subtract their linear combinations, and divide the remainder by $S$.
The resulting series converges uniformly in the compact topology and
defines a continuous surjection $k[[z]]^r\to M/I$.  The kernel and
cokernel of $z$ on a finite module over this discrete valuation ring
are finite-dimensional.  The same conclusion for $S$ on $M$ follows
using its finite invariant subspace $I$.

The image of $\ker f$ in $M/I$ is contained in the $z$-power torsion
submodule, which is finite-dimensional.  Since $I$ is finite, this
proves that $\ker f$ is finite, without assuming it closed in advance.
For $x$ on the left of \eqref{eq:compact:finite-intersection}, its fibre
is now a finite closed coset of $\ker f$.  It meets every nested compact
set $S^aM$, and hence meets $I$.  This proves the equality.
\end{proof}

We will also use that a countable compact Hausdorff group is finite.
The Baire theorem applied to its countably many singleton subsets gives
an isolated point; translations make the group discrete, and
compactness then makes it finite.  As $k$ is finite, countable
dimension and countability of the underlying set agree.

\subsection{The regular translation module}

For $t>1$, the translation representations supply the missing
countability argument. We work simultaneously at sufficiently deep
frame levels and choose the largest degree in which countability
could fail. A Koszul extension will give one surjective operation
into that degree, modulo countable-dimensional spaces. Its powered
realizations must all be compared at one fixed frame level: only then
can the finite transfer intersection in the compact operator lemma
bound its image. The next two lemmas construct this operation; the
following subsection establishes the required uniformity.

Assume first that $t>1$, so that
$\Lambda=k[[D_1,\ldots,D_d]]$ with $d=m(t-1)>0$.
We use the finite-dimensional nilpotent $\Lambda$-modules, equivalently
the finite translation representations.  Associated coefficients,
followed by $P_n$-cochains, form an exact functor on cochain terms:
finite flat torsor descent and the split module filtrations above give
continuous coefficient splittings.  The finite $P_n$-resolution gives
one bound on its cohomological degrees.

To retain precisely the possible failure of countability, let
$\mathscr V$ be the Serre quotient of $k$-vector spaces by the
countable-dimensional ones.  Let $\mathscr G$ be the category of germs
of sequences in $\mathscr V$, where sequences agreeing after an initial
segment are identified.  Use a cofinal sequence of principal joint
frame levels.  Denote by $T^i(N)\in\mathscr G$ the cohomology of the
associated coefficient for a finite module $N$ at those levels.
In particular $T^i(k)$ is represented by $X^i$.
For each finite diagram we may discard an initial segment to make all
its finite auxiliary actions defined and compatible.  Kernels,
cokernels, and exact sequences have their termwise meaning in this
category.  Its zero objects are exactly the sequences of
countable-dimensional spaces at all sufficiently deep levels.
Thus $T^i(k)\ne0$ means that $X^i$ is uncountable-dimensional at
arbitrarily deep levels; discarding a finite initial segment cannot
remove that failure.

The exact cochain functor extends levelwise to pro objects.  Write
$\overline T^i(\Lambda)$ for its value on the inverse regular-module
system.  Lemma~\ref{lem:cartier-transfer} identifies this system with
\[
          \cdots\xrightarrow{S}X^i\xrightarrow{S}X^i
\]
in $\operatorname{Pro}(\mathscr G)$, with the augmentation to
$T^i(k)$; the harmless normalization constants are in $k^\times$.
The completed Koszul resolutions used below are exact pro resolutions.
Indeed finite-length quotients of a finite exact complex of finitely
generated $\Lambda$-modules are pro exact by the adic Artin--Rees
property.  Thus they can be used with this levelwise cochain functor.

\begin{lemma}\label{lem:compact:pro-constant}
For each $i$, the pro object $\overline T^i(\Lambda)$ is constant and
its augmentation to $T^i(k)$ is monic.  Its $\Lambda$-action factors
through $k$.
\end{lemma}

\begin{proof}
Constant objects in a pro abelian category are closed under kernels
and cokernels of maps between constants and under extensions.  For the
last assertion, one can pass to the opposite ind category: in an
extension of two constants, the identity of the constant quotient
factors through one presenting term.  Add the constant kernel to that
term; the resulting relation object is constant, giving a constant
presentation of the extension.

We induct downwards in $i$, starting above the uniform cohomological
bound.  Suppose the assertion of constancy is known for larger indices.
Resolve $\mathfrak m_\Lambda=(D_1,\ldots,D_d)$ by the truncated
completed Koszul resolution.  Successive short exact sequences
$0\to K\to F\to Q\to0$ in this finite resolution give
\[
0\longrightarrow
 \coker\bigl(\overline T^j(K)\to\overline T^j(F)\bigr)
 \longrightarrow\overline T^j(Q)
 \longrightarrow
 \ker\bigl(\overline T^{j+1}(K)\to\overline T^{j+1}(F)\bigr)
 \longrightarrow0.
\]
Starting from its free end, closure of constants under these operations
shows that $\overline T^j(\mathfrak m_\Lambda)$ is constant for
$j>i$.  The augmentation sequence then identifies the image of
$\overline T^i(\Lambda)\to T^i(k)$ with the constant kernel of
$T^i(k)\to\overline T^{i+1}(\mathfrak m_\Lambda)$.

This image is represented by the decreasing system $S^aX^i\subset X^i$.
It stabilizes at a single finite index $h$ in $\mathscr G$.  More
explicitly, after removing its constant subobject the decreasing
system is pro zero.  A sufficiently late transition to any fixed term
is therefore zero; since that transition is monic, its source is zero.
This gives stabilization, and hence
\[
              S^hX^i/S^{h+1}X^i
\]
is countable-dimensional at every sufficiently deep frame level.

At each such level put $M=S^hM^i$.  The map from $M$ to $S^hX^i$ has
countable kernel and cokernel by the strip argument.  It follows that
$M/SM$ is countable, and it is a compact quotient because $SM$ is
compact.  It is therefore finite.  Its stable part is finite by
\eqref{eq:compact:finite-stable-part}.  Lemma~\ref{lem:compact:operator}
now shows that $S$ has finite kernel and cokernel on $M$, and hence
countable kernel and cokernel on $S^hX^i$.  Thus $S$ is an isomorphism
on the stabilized image in $\mathscr G$.  The stationary inverse
system is consequently constant, and its augmentation is the inclusion
of that stabilized image in $T^i(k)$, which is monic.  This completes
the induction.  Finally the augmentation is $\Lambda$-linear and
$\mathfrak m_\Lambda$ acts trivially on $k$; its monicity implies the
same assertion on $\overline T^i(\Lambda)$.
\end{proof}

Suppose now, for a contradiction, that some $T^i(k)$ is nonzero, and
choose the largest such degree $b_0$.  Every finite nilpotent
$\Lambda$-module has a filtration by copies of $k$, so its $T^j$
vanishes for $j>b_0$.  The augmentation sequence in the preceding
proof is therefore surjective in degree $b_0$ as well as injective.
The stabilized image is already the entire $T^{b_0}(k)$, and one may
take $h=0$.  In particular $M^{b_0}/SM^{b_0}$ is finite.  The
locally $S$-nilpotent kernel $M^{b_0}\to X^{b_0}$ is finite by
Lemma~\ref{lem:compact:operator}, and at every sufficiently deep level
\begin{equation}\label{eq:compact:top-finite-intersection}
 I_X=\bigcap_{a\geq0}
       \operatorname{im}(S^aM^{b_0}\longrightarrow X^{b_0})
 \quad\text{is finite}.
\end{equation}

Here $P$ is the fixed target lattice defining $M^{b_0}$. We will
construct a connecting operation into $X^{b_0}$ whose image, on a
sufficiently positive \emph{source} lattice, lies in $I_X$. The source
lattice will be in an associated translation coefficient, and need
not be $P$.

\begin{lemma}[A fixed Koszul edge]\label{lem:compact:koszul-edge}
For one sufficiently large allowed exponent $s$, a generator
$e_s\in\Ext^d_\Lambda(N_s,k)$ induces an epimorphism
\[
             T^{b_0-d}(N_s)\longrightarrow T^{b_0}(k)
\]
in $\mathscr G$.
\end{lemma}

\begin{proof}
Use the completed free Koszul complex
$K_s=K(D_1^s,\ldots,D_d^s)$ in columns $[-d,0]$.
Here $s$ is an allowed distribution exponent from
Lemma~\ref{rings:dual-distribution}, not a height.
Its cohomology is $N_s=\Lambda/(D_1^s,\ldots,D_d^s)$ in column zero.
After applying the cochain functor the vertical page consists of
copies of $\overline T^j(\Lambda)$.  The first differential is zero
by Lemma~\ref{lem:compact:pro-constant}.
For $s'>s$, the map lifting the upward inclusion of regular modules is
given on the Koszul basis by
\[
 e_I\longmapsto
       \left(\prod_{j\notin I}D_j^{s'-s}\right)e_I,
                 \qquad |I|=-\text{column}.
\]
It is the identity in column $-d$ and is
$\mathfrak m_\Lambda$-valued in every other column.  On vertical
cohomology the comparison is therefore the identity in the leftmost
column and zero elsewhere.

There are no incoming differentials to the entry $(-d,b_0)$.
Increase $s$ once for each possible outgoing differential.  Naturality
and the zero map on its target show successively that it vanishes at
the increased stage.  Only finitely many increases are needed, since
there are $d+1$ columns.  The whole leftmost entry therefore survives
and is the quotient edge of $T^{b_0-d}(N_s)$.  Projection to the
leftmost column followed by augmentation represents a nonzero generator
of $\Ext^d_\Lambda(N_s,k)$, proving the assertion.

The operation just described is also the operation of a finite Yoneda
extension.  Here are the category comparisons needed for this use.
Pro finite-length $\Lambda$-modules are the corresponding profinite
modules; stable images identify their finite quotients.  The completed
free modules in $K_s$ are projective in these resolutions.  Moreover
Ext between finite modules over the Noetherian local ring $\Lambda$
agrees with Ext in its $\mathfrak m_\Lambda$-torsion category, since
injective envelopes of torsion modules remain torsion.  A Yoneda
extension in that category can be replaced by a finite one: lift a
finite set of generators successively through its finite number of
arrows and include the resulting finite-length submodules and their
relations.  The same construction contains any prescribed finite
comparison diagram.  Hence completed resolutions, torsion-module Ext,
and finite Yoneda diagrams induce the same connecting operation under
our exact cochain functor.
\end{proof}

\subsection{Uniform equivariance of the extension operation}

The Koszul edge has fixed its endpoints $N_s$ and $k$. We now compare
it with operations obtained at growing root levels $q$. Equality in
ordinary translation Ext would allow a different auxiliary subgroup
for each $q$. We need one subgroup on which all these equalities hold,
because $I_X$ is an intersection at one frame level.

\begin{lemma}[Uniform restriction for fixed endpoints]
\label{lem:uniform-ext}
Choose finite extension diagrams for $e_s$ as above and for a nonzero
$e\in\Ext^d_\Lambda(k,k)$.  They are equivariant for some compact
open frame group $U$.  For every allowed $q\geq s$ let $\gamma_q$ be
the composite
\begin{equation}\label{eq:compact:coinduced-extension}
 N_s\longrightarrow N_q
 \xrightarrow{\operatorname{coind}(e)}N_q[d]
 \longrightarrow k[d],
\end{equation}
where the outer maps are upward inclusion and normalized transfer.
There is one open subgroup $U'\subset U$, independent of $q$, on
which
\[
                         \gamma_q=c_qe_s,
                       \qquad c_q\in k^\times,
\]
as equivariant extensions.
\end{lemma}

\begin{proof}
Forget the auxiliary group first.  Finite Hopf duality identifies
coinduction with base change $D_i\mapsto D_i^q$ of the fundamental
Koszul class.  Its top coefficient is a unit.  The comparison from
$K_s$ to $K_q$ is the identity on its leftmost column, as in
Lemma~\ref{lem:compact:koszul-edge}.  Hence the composite is a nonzero
multiple $c_qe_s$ in ordinary $\Lambda$-Ext.

We must make all these equalities equivariant at a single level.
Their differences have the fixed endpoints $N_s,k$, so they lie in
the kernel of the one forgetful map
\[
 \Ext^d_{T\rtimes U}(N_s,k)
       \longrightarrow\Ext^d_\Lambda(N_s,k).
\]
We will show that this kernel is finite and that each of its classes
vanishes on an open subgroup. A finite intersection will then work
for every $q$.
Shrink $U$ to a pro-$p$ compact analytic subgroup defining the two
chosen finite diagrams.  There is a continuous spectral sequence
\[
 H^a\bigl(U,\Ext^j_\Lambda(N_s,k)\bigr)
 \Longrightarrow \Ext^{a+j}_{T\rtimes U}(N_s,k).
\]
It can be formed in rational translation representations with smooth
auxiliary action.  To verify the needed acyclicity, use cofree
comodules for the semidirect product: their restriction to $T$ is
cofree, and their $T$-invariants are cofree for $U$.  This gives the
usual composition-of-invariants spectral sequence.  The Koszul
resolution shows
\[
                  \dim_k\Ext^j_\Lambda(N_s,k)=\binom dj.
\]
Its terms have continuous finite-dimensional $U$-action.  Analytic
group cohomology is finite-dimensional on such coefficients, so
$\Ext^d_{T\rtimes U}(N_s,k)$ is finite.  Since $k$ is finite, it is a
finite set as well.

The finite Yoneda construction also gives finite equivariant
representatives here. A finite set of vectors has finite $U$-orbits,
since the action is smooth and $U$ is compact. Their $\Lambda$-span
is finite-dimensional: one power of the augmentation ideal kills
the finite orbit set, and that ideal is $U$-stable. Using these spans
successively in the construction retains equivariance.

Every class in the forgetful kernel vanishes after restriction to
some open subgroup. Indeed choose a finite Yoneda zero-comparison for its
underlying extension, using the finite comparison property in the
preceding proof.  All its objects are killed by one power of
$\mathfrak m_\Lambda$. Shrink the auxiliary group so that it fixes
that finite quotient of $\Lambda$ and the finite objects in the
original diagram. Give the additional comparison objects the trivial
action of this subgroup. This is compatible with their $\Lambda$-action,
and all the comparison maps are now equivariant. Intersect these open subgroups
over the finite forgetful kernel.  The resulting $U'$ kills that
entire kernel at once, and hence kills every
$\gamma_q-c_qe_s$.

All the composites in \eqref{eq:compact:coinduced-extension} are
defined before this common restriction by natural coinduction and
transfer.  In particular the transfer $N_q\to k$ is equivariant; its
translation-Hom space is one-dimensional, and a pro-$p$ group has no
nontrivial character into $k^\times$.  The argument uses the fixed
endpoints $N_s,k$.  It requires no common trivialization of the
growing representations $N_q$.
\end{proof}

\subsection{Countability and compact finiteness}

We finish the contradiction following
\eqref{eq:compact:top-finite-intersection}. Fix the finite extension
diagrams first, and then fix a frame level deep enough for
Lemma~\ref{lem:uniform-ext}, the Koszul epimorphism, and all the
preceding compact estimates. At this level choose a free source
lattice $L_s$ in the coefficient associated to $N_s$. For an
arbitrarily deep positive multiple $x^aL_s$, the
cokernel on localized cohomology is countable-dimensional by the strip
argument.  Thus, modulo a countable-dimensional contribution, source
classes are represented by cocycles in one such positive lattice.

Keep separate the two exponents from
Lemma~\ref{rings:dual-distribution}. At a growing divisible torsion
level $\ell_r$ they are
\[
 Q_r=p^{t\ell_r},\qquad q_r=p^{t\ell_r/(t-1)}.
\]
The module $N_{q_r}$ has distribution exponent $q_r$, whereas powering
its coefficient functions uses $Q_r$. Take $q_r\geq s$.
There is a fixed loss $c_{\mathrm{in}}$ in the source-basis and
root-inclusion comparison before powering. Indeed the finitely many
basis functions of $L_s$ belong to one root ring, and their powers in
the base ring have a common denominator. There is also a fixed loss
$c_{\mathrm{out}}$ for the connecting operation of $e$ after powering:
construct it with continuous module sections in its finitely many
finite free coefficient modules. Their section matrices,
differentials, and multiplication structure constants have bounded
denominators. These two constants absorb all changes among the fixed
finite bases. Choose $a>c_{\mathrm{in}}$ once, before varying $r$.
The resulting lower bound for the order is
\[
                  Q_r(a-c_{\mathrm{in}})-c_{\mathrm{out}}.
\]
The first loss is multiplied by the Frobenius exponent, and the
second is incurred by the fixed boundary diagram afterwards. The
bound tends to infinity. Lemma~\ref{lem:cartier-transfer} applies to
the entire finite diagram, so its coinduced connecting operation is
exactly this powered operation on the root coefficient, with all
auxiliary actions retained.

Consequently, for every sufficiently large allowed $q_r$, the powered
representative of a class from the chosen positive lattice is carried
by the fixed boundary of $e$ into $P$ at that root stage.  The final
transfer in \eqref{eq:compact:coinduced-extension}, in stationary
coordinates, is an arbitrarily high iterate of $S$ as $r$ grows.
Let $e_{s,*}$ and $\gamma_{q_r,*}$ denote the connecting operations on
cohomology. For a source class $[z]$ represented in $x^aL_s$ and any
$a_0$, choosing $r$ sufficiently large gives
\[
 c_{q_r}e_{s,*}([z])=\gamma_{q_r,*}([z])
       \in\operatorname{im}(S^{a_0}M^{b_0}\to X^{b_0}).
\]
The equality holds at our single fixed frame level by
Lemma~\ref{lem:uniform-ext}. Its scalar is nonzero, so the
$e_s$-image belongs to all the displayed subspaces. It lies in the
finite space $I_X$ of \eqref{eq:compact:top-finite-intersection}.

Thus the image of the whole Koszul edge is countable-dimensional:
its positive-lattice part is finite and its remaining source quotient
is countable-dimensional.  The edge was an epimorphism modulo
countable-dimensional spaces at every sufficiently deep level, so
$X^{b_0}$ is countable-dimensional at all those levels.  This says
$T^{b_0}(k)=0$, a contradiction.  We have proved that every $X^i$
is countable-dimensional at sufficiently deep joint frame levels.

At $t=1$ the translation group is diagonalizable and one replaces this
pro argument by its exact character decomposition.  Constants split
equivariantly at every root step, and the normalized constant-character
projection retracts inclusion. Fix a divisible stationary step with
function exponent $Q$ fixing $k$. In stationary coordinates the
retraction says
\[
 S^r\Phi_{Q^r}=1,\qquad \Phi_{Q^r}(f)=f^{Q^r},
\]
as maps of coefficient complexes. Apply the same
positive-lattice estimate: a class represented in one sufficiently
positive lattice is, after powering and then transfer, represented in
$S^aM^i$ for arbitrarily large $a$.  The retraction identifies its
class with the original one already in compact lattice cohomology.
It therefore belongs to the finite intersection
$\bigcap_aS^aM^i$ from
\eqref{eq:compact:finite-stable-part}.  Pole strips account for a
countable-dimensional remainder.  This proves the same countability
of $X^i$, using neither a positive-dimensional distribution ring nor
higher translation Ext in multiplicative height.

\begin{proposition}[Compact coefficient finiteness]
\label{prop:compact-finiteness}
The first assertion of Theorem~\ref{thm:coefficient-finiteness} holds
at $(n,t)$ under the earlier compact assertions in
\eqref{eq:compact:induction-order}.  In addition, all finite-frame
localized coefficients, and the localizations of admissible compact
coefficients, have countable-dimensional $P_n$-cohomology.
\end{proposition}

\begin{proof}
We have proved countability for functions at sufficiently deep finite
frame levels.  It descends to any smaller level by finite flat
\v{C}ech descent.  On cochain terms augmentation is exact: finite
\'{e}tale descent has module sections, or a trace-one contraction,
and these preserve bounded poles.  Pass first to a normal finite frame cover, with deck group $D$.
Its \v{C}ech coefficient in degree $q$ is a finite product of the
deeper-level coefficient indexed by $D^q$.  The acting stabilizer
$P_n$ fixes this indexing set because its action commutes with
the frame action.  Thus each such term is a finite direct sum
of the already treated $P_n$-coefficient; no new acting-stabilizer
representation is introduced here.  In each total degree only finitely
many terms of its bounded-below cohomological spectral sequence
contribute.  Countability is therefore preserved in the descent.

For an admissible extra coefficient, take the finite frame and
filtration in Definition~\ref{def:compact:admissible}.  Trivial
graded pieces reduce to the function calculation.  If a finite
translation splitting is required, filter its representation over
the marked unsplit base.  At height greater than one the graded
pieces are trivial.  At height one a character is a summand of the
regular representation of the finite diagonalizable quotient; its
associated coefficient is a summand of a finite root ring.  Powering
identifies this ring, with all actions transported, with the function
ring at an adequately deep marking level. The established
countability applies to it and its summands. Enlarge any prescribed
root ascent to one of the cofinal divisible lift levels; the split
filtration pulls back to this level. After the finite marking, its
descent uses the translation torsor of Proposition~\ref{prop:frame-torsor}.
Each degree of its bar construction adds a
finite tensor power of the Honda torsion-coordinate algebra. Its
$P_n$-action is trivial, so the resulting coefficients are finite
direct sums of the treated coefficient. The underlying module
splittings preserve bounded poles, and only finitely many bar terms
contribute in a fixed total degree. Finite extensions and these two
descents prove countability for the localized extra.

Now let $E$ be a compact admissible coefficient on $C_1$.
Its map to $E[1/x]$ has countable-dimensional cohomological kernel by
the strip argument: the boundary restrictions and their conormal
twists are admissible and have finite cohomology by induction.
Its localized cohomology is countable-dimensional by the preceding
paragraph.  Hence $H^i(P_n,E)$ is countable-dimensional.  It is also
profinite by Lemma~\ref{lem:compact:duality}, and is therefore finite.
Only finitely many degrees occur.

For $E$ on $A$ first localize at $x_t$ and pass to the first labeled
level.  The exact-height countability just proved descends to the
unlabeled localization.  The strips for powers of $x_t$ have a finite
filtration by coefficients on the next starting-height disc, with
the specified conormal twists.  These satisfy the induction
hypothesis after pullback to its labeled strata.  The identical
countable-kernel argument and compactness prove finiteness on $A$.

The argument is unaffected by a fixed finite root change: transport
the group, its marking and splitting torsors, the filtrations, and all
twists by that same power identification.  In particular a universal
finite torsion construction restricted to a dead-label boundary may
be a Frobenius pullback of the smaller universal construction; its
frames and splittings are transported from those of that construction.
This checks exactly the extra coefficients and their linear duals
used in the smaller-height local duality of
Proposition~\ref{prop:pole-removal}.
\end{proof}

The distinction between the two conclusions here is essential for the
next step: compact cohomology is finite, whereas the argument so far
gives only countability for the algebraic bounded-pole localization.
Finiteness of that localization will follow from its boundary support
rows in the next section.
